\documentclass[11pt]{article}
\usepackage{amsmath,amssymb,amsthm,mathtools}
\usepackage{enumitem}
\usepackage[margin=1in]{geometry}

\newtheorem{theorem}{Theorem}
\newtheorem{definition}{Definition}
\newtheorem{nogo}{No-Go Theorem}

\title{Step 161: Boundary-Symbol Theorem Extraction / Scope Closure}
\author{Six Birds / RH Membrane Working Note}
\date{}

\begin{document}
\maketitle

\section{Orientation and inherited state}

This step is \emph{adequacy-oriented}.  Smaller residual is better:
\[
  \Xi^{\rm BC}\to 0
  \quad\text{or}\quad
  \Xi^{\rm BC}\preceq \Omega
\]
with a declared budget.

The input state is Step 160.  The active residual is
\[
  \Xi^{\rm BC}_{\ell,a}
  =
  \mathfrak B_{\ell,a}^{*}\Pi_{Y_a}\mathfrak B_{\ell,a},
\]
and the active Calkin object is
\[
  C_\ell P_\eta,
  \qquad
  C_\ell=(I-P_\infty)M_{m_\ell}P_\infty,
  \qquad
  m_\ell(s)=e^{-\ell(1/2-s)}.
\]
Here
\[
  H_\eta=
  \overline{\operatorname{span}\{J_a^*Y^a_{\rho,k}\}},
  \qquad
  P_\eta=\operatorname{Proj}_{H_\eta}.
\]
The comparison context is the Burnol/Sonine pulled-evaluator carrier in the
Calkin algebra \(\mathcal B(H)/\mathcal K(H)\).  Hard-support and
Hardy/Hankel computations are not in this context unless a bridge is supplied.

\section{Bridge theorem target}

\begin{definition}[Calkin-faithful boundary-symbol theorem]
A Calkin-faithful boundary-symbol theorem for the Step 161 target consists of:
\begin{enumerate}[label=(\roman*)]
  \item a declared Calkin algebra \(\mathcal A_\eta/\mathcal K_\eta\) generated
  by the Burnol/Sonine operators \(P_\infty\), \(M_{m_\ell}\), and \(P_\eta\);
  \item a quotient-level boundary-symbol map
  \[
    \sigma_B:\mathcal A_\eta/\mathcal K_\eta\to\mathcal S_B
  \]
  that is faithful on the Calkin class of \(C_\ell P_\eta\);
  \item a factorization on the completed Burnol/Sonine carrier
  \[
    C_\ell=U_\ell W_\ell+K_\ell,
    \qquad
    K_\ell\in\mathcal K(H);
  \]
  \item lower-faithfulness of \(U_\ell\) on the boundary range relevant to
  \(H_\eta\);
  \item compact-ideal compatibility, so that compactness of \(W_\ell P_\eta\)
  is equivalent to compactness of \(C_\ell P_\eta\).
\end{enumerate}
\end{definition}

This theorem is not a notation choice.  It is a completed-carrier operator
theorem.  A source that supplies only a Hardy shadow, hard-support shadow,
finite singular-value pattern, or ambient Sonin formula does not satisfy the
target unless it also supplies the bridge above.

\section{Source audit}

Step 161 audits the local sources available after Step 160.

\begin{center}
\begin{tabular}{lll}
\hline
Source & Status & Reason\\
\hline
Burnol carrier theorem & accepted support & defines \(Y_a\), \(P_a=Y_a^\perp\), and evaluators\\
CCM/Sonin machinery & imported support & supplies semilocal/Sonin setting, not Calkin bridge\\
Hardy/Hankel shadow & public shadow & wrong carrier without bridge\\
Hard-support shadow & public shadow & gives obstruction model, not Burnol theorem\\
Step 160 theorem & conditional & bridge equivalence only if bridge exists\\
\hline
\end{tabular}
\end{center}

No audited local source supplies all bridge gates.  Therefore the bridge is not
accepted in this step.

\section{Conditional theorem retained}

\begin{theorem}[Conditional boundary-symbol equivalence]
Assume the Calkin-faithful boundary-symbol theorem above.  Then
\[
  C_\ell P_\eta\in\mathcal K(H)
  \quad\Longleftrightarrow\quad
  W_\ell P_\eta\in\mathcal K(H).
\]
Consequently the compact/tail status of \(\Xi^{\rm BC}\) is equivalent, under
the bridge, to the compact/tail status of
\[
  \Xi^{\rm BW}_{\ell,a}
  =
  P_\eta W_\ell^*W_\ell P_\eta.
\]
\end{theorem}

\begin{proof}[Proof sketch]
The factorization \(C_\ell=U_\ell W_\ell+K_\ell\) gives
\[
  C_\ell P_\eta=U_\ell W_\ell P_\eta+K_\ell P_\eta.
\]
The second term is compact.  If \(W_\ell P_\eta\) is compact, then
\(C_\ell P_\eta\) is compact.  Conversely, after passing to the Calkin algebra,
faithfulness of \(\sigma_B\) and lower-faithfulness of \(U_\ell\) on the
boundary range prevent \(U_\ell\) from killing a noncompact boundary class.
Thus a noncompact class for \(W_\ell P_\eta\) would give a noncompact class for
\(C_\ell P_\eta\), proving the equivalence.
\end{proof}

\section{No-go records retained}

\begin{nogo}[Public-shadow non-promotion]
Without a Calkin-faithful boundary-symbol bridge, Hardy/Hankel and
hard-support compactness or noncompactness statements cannot be promoted to
the Burnol/Sonine pulled-evaluator carrier.
\end{nogo}

\begin{nogo}[Finite-window Calkin blindness]
Finite-section singular-value or matrix evidence cannot decide whether
\[
  C_\ell P_\eta\in\mathcal K(H)
\]
or whether \(\|C_\ell P_\eta\|_{\rm e}>0\) on the completed carrier.
\end{nogo}

Both no-go statements are framework-derived analytical structural content.
They do not rule out a true bridge, a direct pulled-evaluator Weyl sequence, a
co-Poisson factorization, or a semilocal prolate cancellation theorem.

\section{Verdict}

\[
\boxed{\text{No Calkin-faithful boundary-symbol theorem is extracted locally.}}
\]

The step is diagnostic-complete at this leaf.  The exact remaining external
theorem is:
\[
  \boxed{
  \text{Construct }\sigma_B,\ U_\ell,\ W_\ell,\ K_\ell
  \text{ with }
  C_\ell=U_\ell W_\ell+K_\ell,\ K_\ell\in\mathcal K,
  }
\]
on the Burnol/Sonine pulled-evaluator carrier, with lower-faithfulness and
compact-ideal compatibility on \(H_\eta\).

\section{Next leaf}

Step 162 should either:
\begin{enumerate}[label=(\alph*)]
  \item attempt the external operator theorem above from Burnol/Sonine or CCM
  source material;
  \item construct an actual pulled-evaluator Weyl sequence in \(H_\eta\);
  \item prove a genuine co-Poisson factorization for the shifted boundary
  packets; or
  \item close the RH membrane packet as a scoped residual theorem carrying
  \(\Xi^{\rm BC}\).
\end{enumerate}

\end{document}
